
\section{Conclusioni}%
\label{sec:conclusions}

\subsection{Risultati e limitazioni}

Il progetto ha portato alla realizzazione di un sistema distribuito per il monitoraggio del livello dell'acqua e il controllo dell'apertura di una valvola.
L'architettura adottata affida al CUS la gestione centralizzata dello stato e della politica di controllo, integrando acquisizione delle misure, attuazione e interfacce operatore locali e remote.

Le principali limitazioni riguardano l'assenza di un feedback sulla posizione effettiva della valvola e il rilevamento delle interruzioni della comunicazione seriale.

\subsection{Possibili sviluppi futuri}

Tra i possibili miglioramenti rientrano l'introduzione di un feedback sulla posizione della valvola, il monitoraggio della comunicazione tra CUS e WCS, la persistenza dello storico delle misure e l'autenticazione delle richieste di controllo.

Queste estensioni potrebbero migliorare l'affidabilità e la gestione del sistema, pur non rientrando negli obiettivi dell'attuale implementazione.
